\section{Conclusion}
\label{sec:conclusion}

This paper constructed four things: an information-safe simulator for an
explicitly specified dialect of six-player Canadian Fish, each consequential
rule a switch and each agent's view audited against the public record
(\S\ref{sec:engine}); an exact observer-conditioned deal posterior, closed-form
because card movement in Fish is public so the hidden state is exactly the
initial deal, giving ownership marginals, team-ownership and named-allocation
probabilities and validated against a brute-force enumeration oracle on small
reachable states (\S\ref{sec:belief}, \S\ref{sec:results-verify}); the cheaper
approximation of it that the deployed \vfast{} policy runs, beneath fitted ask,
lookahead and declaration machinery (\S\ref{sec:policy}); and a duplicate-block
protocol --- every primary-bank deal in six seat rotations, two on the ablation
bank, intervals bootstrapped with the deal as the resampling cluster
(\S\ref{sec:protocol}).

Four findings are empirical. \vfast{} defeated \vthree{} in \numVsVthreeWins{}
of \numHtHGames{} games --- \numHtHDeals{} held-out deals in six seat rotations
--- at \numVsVthree\%, deal-clustered 95\% interval \numVsVthreeCI\%: its lowest
rate against any member of this eight-policy panel, on this bank, under this
dialect and the default arbitration order, six of the eight also being
fitting-panel members (\S\ref{sec:results-h2h}). The margin is declaration
reliability rather than ask conversion: on the same games \vfast{} declared
correctly \numVsVthreeDeclAcc\% of the time against \numVthreeDeclAccSame\% for
\vthree{}, at expected calibration error \numDeclECE{} against \numVthreeDeclECE{}
(\S\ref{sec:results-calib}); no experiment decomposes that margin causally.
Among the frozen-policy ablations, the rows removing parts of the rules-derived
constraint set, the policy prior, the lock-completion weight and the \vblock{}
belief substitution are resolved from zero; the declaration timing rule, the
two-ply refinement and the information-leak weights are not, and support no
conclusion in either direction (Table~\ref{tab:ablations}). And the two cheap
declaration pre-gates are a pruning heuristic, not a proved bound: over the
primary evaluation corpus they rejected \numGateFalseNeg{} candidate pairs the
ungated rule would have declared --- a measured, non-zero false-negative rate
(E16, \S\ref{sec:results-verify}).

Three readings are suggested by the analysis without being established by it.
The declaration-accuracy difference appears carried by the joint allocation
probability rather than by the timing rule above it: a fixed threshold in place
of the fitted stopping rule moved the win rate \numAblStop{} points, 95\% paired
interval \numAblStopCI{} (E5) --- an interval that resolves nothing, beside a
large and consistent reliability difference. Substituting the \vblock{} belief
cost \numAblBlock{} points, 95\% paired interval \numAblBlockCI{}, so a policy
fitted around one belief representation does not transfer to the other
unchanged; that does not show an exact-belief policy would be weaker after
matched retraining, and no matched fit was run. And ownership of a half-suit can
freeze while allocation information about it keeps arriving, asks inside a
locked half-suit remaining legal for the owning team and still carrying
certificates (\S\ref{sec:locked}): waiting is safe with respect to ownership
without being profitable, and this study does not isolate its value.

More remains unresolved than settled. Exploitability is not bounded: the
response probe returns a lower bound within one parametric class, its interval
including even money (\S\ref{sec:lbr}). No team equilibrium was computed.
Termination is imposed by hand-set horizons rather than derived, and the
observed cycling is a property of specific fitted configurations in self-play,
not of the game. Nothing was tested against expert human play or an
independently authored engine. The open question the belief-substitution and
response-probe results jointly raise is whether a separately fitted
exact-belief agent, a policy class broader than a linear ask score with a linear
value function, or an equilibrium-oriented method can improve on the current
\vfast{} policy.
